\section{Simple differential equations}

\begin{definitionbox}
A differential equation is an equation involving an unknown function and one or more of its derivatives.
\end{definitionbox}

Such an equation describes a relation between a function and its change. Unlike an algebraic equation, it does not usually ask for a single number. It asks for a function.

The equation
\[
\frac{dy}{dx}=3
\]
asks for a function whose derivative is \(3\). The general solution is
\[
y(x)=3x+C.
\]

Another simple example is
\[
\frac{dy}{dx}=2x.
\]
A solution is
\[
y(x)=x^2+C,
\]
because
\[
\frac{d}{dx}(x^2+C)=2x.
\]

A slightly different example is
\[
\frac{dy}{dx}=ky,
\]
where \(k\) is a constant. This equation says that the derivative of the function is proportional to the function itself. One family of solutions is
\[
y(x)=Ce^{kx}.
\]
Indeed,
\[
\frac{d}{dx}Ce^{kx}=kCe^{kx}=ky(x).
\]

The constant \(C\) is determined by an additional condition. For example, if
\[
y(0)=2,
\]
then
\[
2=Ce^{k\cdot 0}=C,
\]
so
\[
y(x)=2e^{kx}.
\]

\begin{remarkbox}
Here the important point is the basic structure: a differential equation contains derivatives and asks for an unknown function.
\end{remarkbox}
